\section{Methodology}
\subsection{Stage 1: BM25 Retrieval}
Given a query $q$ (the candidate-profile skill tokens) and a document $d$
(a job posting's title, skills, seniority, and location), BM25 scores
\[
\mathrm{BM25}(q,d)=\sum_{t\in q} \mathrm{IDF}(t)\,
\frac{f(t,d)\,(k_1+1)}{f(t,d)+k_1\big(1-b+b\,\tfrac{|d|}{\mathrm{avgdl}}\big)},
\]
where $f(t,d)$ is term $t$'s frequency in $d$, $|d|$ is document length,
$\mathrm{avgdl}$ is the mean document length over the indexed corpus, and
$k_1$, $b$ are the usual saturation and length-normalization constants. The
inverse document frequency term is
\[
\mathrm{IDF}(t)=\ln\!\left(\frac{N-n(t)+0.5}{n(t)+0.5}+1\right),
\]
with $N$ indexed documents and $n(t)$ the number containing $t$. The index
is fit on exactly the training partition's candidate pool (\texttt{rank\_bm25}),
so stage-1 retrieval itself obeys the \texttt{fit\_on: train\_only} contract.

\subsection{The Taxonomy-Normalized Overlap Feature}
BM25 rewards surface-form overlap: a query token must match a document
token exactly. Skill mentions in practice are not canonical ("JS",
"Javascript", "ECMAScript" are one skill under three spellings), so a
taxonomy-normalized Jaccard overlap is computed separately and fed to
stage~2 as a feature:
\[
\mathrm{overlap}(C,J) = \frac{|\mathrm{skills}(C)\cap\mathrm{skills}(J)|}
{|\mathrm{skills}(C)\cup\mathrm{skills}(J)|},
\]
where $\mathrm{skills}(\cdot)$ denotes the canonical (synonym-collapsed)
skill set of a candidate profile $C$ or job posting $J$. This is the one
feature a purely lexical retriever cannot express, and it is the single
strongest stage-2 signal by gain-based importance.

\subsection{Stage 2: LambdaMART Re-ranking}
LambdaMART optimizes a pairwise surrogate for a listwise ranking metric.
For a pair of documents $i$ (more relevant) and $j$ (less relevant) within
the same query group, the gradient of the loss with respect to the model's
score $s_i$ is
\[
\lambda_{ij} = -\frac{\partial \mathcal{L}}{\partial s_i}
= \sigma\big(s_j - s_i\big)\cdot|\Delta\mathrm{NDCG}_{ij}|,
\]
where $\sigma(\cdot)$ is the logistic function and $|\Delta\mathrm{NDCG}_{ij}|$
is the change in NDCG that would result from swapping $i$ and $j$ in the
current ranking. Scaling the pairwise gradient by this swap cost is what
makes LambdaMART a listwise approximation rather than a plain pairwise
classifier: mis-ordering two documents near the top of the ranked list is
penalized far more than mis-ordering two documents far down it, which
directly targets the NDCG objective stage~2 is evaluated against. LightGBM's
\texttt{objective=lambdarank} trains exactly this gradient over the
handcrafted feature matrix described in Section~3, using the LTR group
structure from the proxy supervision in Section~2.
